\documentclass[11pt,a4paper,fontset=fandol]{ctexart}

\usepackage[margin=2.3cm]{geometry}
\usepackage{amsmath,amssymb,bm}
\usepackage{booktabs}
\usepackage{array}
\usepackage{longtable}
\usepackage{enumitem}
\usepackage{xcolor}
\usepackage{hyperref}
\usepackage{listings}
\usepackage{microtype}

\hypersetup{
    colorlinks=true,
    linkcolor=blue,
    urlcolor=blue,
    citecolor=blue
}

\lstset{
    basicstyle=\ttfamily\small,
    breaklines=true,
    frame=single,
    columns=fullflexible,
    keepspaces=true
}

\setlength{\parindent}{2em}
\setlength{\parskip}{0.4em}
\renewcommand{\arraystretch}{1.25}

\title{On Policy Distillation Notes}
\author{}
\date{}

\begin{document}

\maketitle
\tableofcontents
\newpage

\section{符号定义}

\begin{itemize}[leftmargin=2em]
    \item \(x=(x_1,\ldots,x_T)\)：学生生成的完整序列。
    \item \(x_{1..t}\)：生成到第 \(t\) 个 token 时的前缀。
    \item \(x_{t+1}\)：下一 token。
    \item \(\pi_\theta\)：学生模型。
    \item \(\pi_{\text{teacher}}\)：固定的教师模型。
    \item \(\theta\)：学生模型参数。
    \item \(\gamma\)：未来 KL 的折扣因子。
    \item \(\beta\)：代码中的 \texttt{kl\_penalty\_coef}。
\end{itemize}

此外还涉及两个新符号用于重要性采样：

\begin{itemize}[leftmargin=2em]
    \item \(q\)：真正生成 rollout 的 sampler policy。
    \item \(p_\theta\)：执行训练 forward 时的 learner policy。
\end{itemize}

理想的严格 on-policy 情况下两者跟 \(\pi_\theta\) 应该都相同；但在实际系统中，rollout 与训练之间通常仍会存在偏差，例如训推Infra等不同，以及复用同一批 rollout 数据进行多次参数更新时，learner 已逐步偏离生成该数据的 sampler。此外，现在有不少infra会通过重要型采样来传 RL 梯度。
\section{从每个 token 的反向 KL 开始}

定义固定前缀 \(x_{1..t}\) 下，学生到教师的下一 token 反向 KL：

\[
\boxed{
\mathrm{KL}\left(\pi_\theta\|\pi_{\text{teacher}}\right)
=
\mathbb E_{x\sim\pi_\theta}
\left[
\log \pi_\theta(x_{t+1}|x_{1..t})
-
\log \pi_{\text{teacher}}(x_{t+1}|x_{1..t})
\right]
}
\]

更严格地展开为全词表求和：

\[
\begin{aligned}
D_t(x_{1..t})
&=
\sum_{x_{t+1}}
\pi_\theta(x_{t+1}|x_{1..t})\\
&\qquad\cdot
\left[
\log\pi_\theta(x_{t+1}|x_{1..t})
-
\log\pi_{\text{teacher}}(x_{t+1}|x_{1..t})
\right].
\end{aligned}
\]

它比较的是：

\[
\pi_\theta(\cdot|x_{1..t})
\quad\text{与}\quad
\pi_{\text{teacher}}(\cdot|x_{1..t}),
\]

即同一个学生访问到的前缀下，两个下一 token 条件分布的反向 KL。

为简化记号，定义单 token 的采样 KL：

\[
d_t
=
\log\pi_\theta(x_{t+1}|x_{1..t})
-
\log\pi_{\text{teacher}}(x_{t+1}|x_{1..t}).
\]

\section{全词表反向 KL 如何得到最终梯度}

固定 \(x_{1..t}\)，对 \(D_t\) 求梯度：

\[
\begin{aligned}
\nabla_\theta D_t
&=
\sum_{x_{t+1}}
\nabla_\theta\pi_\theta(x_{t+1}|x_{1..t})\,d_t\\
&\quad+
\sum_{x_{t+1}}
\pi_\theta(x_{t+1}|x_{1..t})
\nabla_\theta\log\pi_\theta(x_{t+1}|x_{1..t}).
\end{aligned}
\]

使用对数求导：

\[
\nabla_\theta\pi_\theta
=
\pi_\theta\nabla_\theta\log\pi_\theta,
\]

得到：

\[
\nabla_\theta D_t
=
\mathbb E_{x_{t+1}\sim\pi_\theta}
\left[
(d_t+1)
\nabla_\theta\log\pi_\theta(x_{t+1}|x_{1..t})
\right].
\]

其中：

\[
\begin{aligned}
\mathbb E_{\pi_\theta}
\left[
\nabla_\theta\log\pi_\theta
\right]
&=
\sum_{x_{t+1}}
\pi_\theta(x_{t+1}|x_{1..t})
\nabla_\theta\log\pi_\theta(x_{t+1}|x_{1..t})\\
&=
\sum_{x_{t+1}}
\nabla_\theta\pi_\theta(x_{t+1}|x_{1..t})\\
&=
\nabla_\theta 1\\
&=0.
\end{aligned}
\]

所以 \(+1\) 在期望中消失：

\[
\boxed{
\nabla_\theta D_t
=
\mathbb E_{x_{t+1}\sim\pi_\theta}
\left[
d_t
\nabla_\theta\log\pi_\theta(x_{t+1}|x_{1..t})
\right]
}
\]

也就是：

\[
\boxed{
\nabla_\theta D_t
=
\mathbb E_{x_{t+1}\sim\pi_\theta}
\left[
\left(
\log\pi_\theta(x_{t+1}|x_{1..t})
-
\log\pi_{\text{teacher}}(x_{t+1}|x_{1..t})
\right)
\nabla_\theta\log\pi_\theta(x_{t+1}|x_{1..t})
\right].
}
\]

这就是 OPD 最核心的 token-level policy-gradient 形式。

\section{sampled tokens 如何估计全词表梯度}

由于直接使用全词表做反向 KL 太贵了，目前也只有 DeepSeek v4 使用，

大家一般会使用 sampled tokens 或者 topk 进行估计。下面先证明 \(K\) 个 sampled tokens 的 Monte Carlo 平均是无偏的，然后再看现在实践中常用的 \(K=1\) 特例。

第三节最后得到的是：

\[
\boxed{
\nabla_\theta D_t
=
\mathbb E_{x_{t+1}\sim\pi_\theta}
\left[
d_t
\nabla_\theta\log\pi_\theta(x_{t+1}|x_{1..t})
\right].
}
\]

把这个期望重新展开成全词表求和：

\[
\begin{aligned}
\nabla_\theta D_t
&=
\sum_{x_{t+1}}
\pi_\theta(x_{t+1}|x_{1..t})\\
&\qquad\cdot
\left[
\left(
\log\pi_\theta(x_{t+1}|x_{1..t})
-
\log\pi_{\text{teacher}}(x_{t+1}|x_{1..t})
\right)
\nabla_\theta\log\pi_\theta(x_{t+1}|x_{1..t})
\right].
\end{aligned}
\]

也就是说，真实梯度需要枚举每个可能的下一 token，并用学生概率

\[
\pi_\theta(x_{t+1}|x_{1..t})
\]

作为权重加权求和。

如果不枚举全词表，而是从学生分布中采样 \(K\) 个 token：

\[
x_{t+1}^{(1)},\ldots,x_{t+1}^{(K)}
\sim
\pi_\theta(\cdot|x_{1..t}),
\]

其中第 \(k\) 个样本对应：

\[
d_t^{(k)}
=
\log\pi_\theta(x_{t+1}^{(k)}|x_{1..t})
-
\log\pi_{\text{teacher}}(x_{t+1}^{(k)}|x_{1..t}).
\]

对第 \(k\) 个 sampled token，可以构造一个单样本梯度估计：

\[
\hat g^{(k)}
=
d_t^{(k)}
\nabla_\theta
\log\pi_\theta(x_{t+1}^{(k)}|x_{1..t}).
\]

然后对 \(K\) 个样本取平均：

\[
\boxed{
\hat g_K
=
\frac1K
\sum_{k=1}^{K}
\hat g^{(k)}
=
\frac1K
\sum_{k=1}^{K}
d_t^{(k)}
\nabla_\theta
\log\pi_\theta(x_{t+1}^{(k)}|x_{1..t}).
}
\]

先看单个 \(\hat g^{(k)}\) 的期望。因为每个 \(x_{t+1}^{(k)}\) 都来自

\[
\pi_\theta(\cdot|x_{1..t}),
\]

所以：

\[
\begin{aligned}
\mathbb E[\hat g^{(k)}]
&=
\sum_{x_{t+1}^{(k)}}
\pi_\theta(x_{t+1}^{(k)}|x_{1..t})\\
&\qquad\cdot
\left[
\left(
\log\pi_\theta(x_{t+1}^{(k)}|x_{1..t})
-
\log\pi_{\text{teacher}}(x_{t+1}^{(k)}|x_{1..t})
\right)
\nabla_\theta
\log\pi_\theta(x_{t+1}^{(k)}|x_{1..t})
\right]\\
&=
\nabla_\theta D_t.
\end{aligned}
\]

因此 \(K\) 个样本平均后的期望为：

\[
\begin{aligned}
\mathbb E[\hat g_K]
&=
\frac1K
\sum_{k=1}^{K}
\mathbb E
\left[
d_t^{(k)}
\nabla_\theta
\log\pi_\theta(x_{t+1}^{(k)}|x_{1..t})
\right]\\
&=
\frac1K
\sum_{k=1}^{K}
\nabla_\theta D_t\\
&=
\nabla_\theta D_t.
\end{aligned}
\]

所以 \(K\) 个 sampled tokens 的平均是全词表梯度的无偏 Monte Carlo 估计。

现在实践里常取 \(K=1\)。此时估计量退化为：

\[
\boxed{
\hat g_1
=
d_t
\nabla_\theta
\log\pi_\theta(x_{t+1}|x_{1..t}),
\qquad
x_{t+1}\sim\pi_\theta.
}
\]

它仍然无偏，但因为只用一个 sampled token，方差比更大的 \(K\) 更高。

我们的上述推导在工程实现上还漏掉了一个很重要的事情：sampled-token estimator 不能直接把 sampled KL 当作普通 loss 反传，而需要 stop-gradient。

实际上我在一开始推导的时候也完全忘掉了，是看到了官方代码才发现有些不对劲儿（

采样之后，

\[
x_{t+1}\sim\pi_\theta,
\]

这个 \(x_{t+1}\) 在自动微分看来只是一个固定 token；采样概率本身没有显式留在计算图里。因此工程实现要复现的不是 \(\nabla_\theta d_t\)，而是上一节推导出的 score-function 梯度估计器：

\[
\hat g_K
=
\frac1K
\sum_{k=1}^{K}
d_t^{(k)}
\nabla_\theta
\log\pi_\theta(x_{t+1}^{(k)}|x_{1..t}).
\]

实现这个估计器的代理损失应写成：

\[
\boxed{
\mathcal L_{\mathrm{MC}}
=
\frac1K
\sum_{k=1}^{K}
\operatorname{sg}\!\left[d_t^{(k)}\right]
\log\pi_\theta(x_{t+1}^{(k)}|x_{1..t})
}
\]

其中 \(\operatorname{sg}\) 表示 stop-gradient。它只把 \(d_t^{(k)}\) 当成数值权重，让自动微分只穿过后面的 score function：

\[
\nabla_\theta \mathcal L_{\mathrm{MC}}
=
\frac1K
\sum_{k=1}^{K}
\operatorname{sg}\!\left[d_t^{(k)}\right]
\nabla_\theta
\log\pi_\theta(x_{t+1}^{(k)}|x_{1..t}).
\]

如果不加 stop-gradient，写成 \(d_t\log\pi_\theta(x_{t+1}|x_{1..t})\)，自动微分会多出：

\[
\log\pi_\theta(x_{t+1}|x_{1..t})\nabla_\theta d_t,
\]

这里和全词表 KL 的区别是：全词表方法没有离散采样，完整目标

\[
\mathcal L_t^{\text{full}}
=
\sum_{x_{t+1}}
\pi_\theta(x_{t+1}|x_{1..t})
\left[
\log\pi_\theta(x_{t+1}|x_{1..t})
-
\log\pi_{\text{teacher}}(x_{t+1}|x_{1..t})
\right]
\]

已经把每个 token 的概率权重 \(\pi_\theta(x_{t+1}|x_{1..t})\) 显式写在求和里，所以可以直接自动微分整个真实目标，不需要 stop-gradient。

\section{sampled-token 与全词表 loss 的区别}

这两个方法估计的是同一个固定前缀条件 KL，但实现方式不同。

\begin{longtable}{>{\raggedright\arraybackslash}p{0.22\textwidth}
                  >{\raggedright\arraybackslash}p{0.34\textwidth}
                  >{\raggedright\arraybackslash}p{0.34\textwidth}}
\toprule
方面 & 全词表 loss & sampled-token loss \\
\midrule
下一 token & 枚举所有词表 token & 只使用实际采样 token \\
目标 & 精确计算条件 KL & Monte Carlo 估计条件 KL 梯度 \\
教师输出 & 完整词表概率或 logits & 被采样 token 的 logprob \\
方差 & 对当前前缀无采样方差 & \(K=1\) 时方差较大 \\
计算量 & 与词表大小相关 & 每个位置只需 sampled token 的教师 logprob \\
自动微分 & 直接对完整 KL 目标反传 & 使用 policy-gradient surrogate \\
stop-grad & 不需要特殊 stop-grad & \(d_t\) 作为 advantage 时需要 stop-grad \\
\bottomrule
\end{longtable}

\section{未来 Token 的影响与折扣因子}

到这里为止，我们一直在固定前缀 \(x_{1..t}\) 下分析下一 token 的条件 KL。但实际上，我们会碰到一个很常见的问题：

当前 token 改变后，后面会访问到的前缀也会改变；这些未来位置上的 KL 要不要归因给当前 token？

如果把整条 rollout 都看作学生模型生成的样本，一个自然的序列级目标是沿轨迹累加每个位置的 KL：

\[
\mathcal J(\theta)
=
\mathbb E_{x\sim\pi_\theta}
\left[
\sum_{t=0}^{T-1}d_t
\right],
\]

其中：

\[
d_t
=
\log\pi_\theta(x_{t+1}|x_{1..t})
-
\log\pi_{\text{teacher}}(x_{t+1}|x_{1..t}).
\]

这个目标也可以理解为完整序列分布上的反向 KL，因为：

\[
\log\pi_\theta(x)
=
\sum_t
\log\pi_\theta(x_{t+1}|x_{1..t}),
\]

教师同理。

对 \(\mathcal J\) 求梯度：

\[
\begin{aligned}
\nabla_\theta\mathcal J
=
\mathbb E_{x\sim\pi_\theta}
\Bigg[
&
\left(
\sum_t\nabla_\theta
\log\pi_\theta(x_{t+1}|x_{1..t})
\right)
\left(
\sum_k d_k
\right)\\
&+
\sum_k
\nabla_\theta
\log\pi_\theta(x_{k+1}|x_{1..k})
\Bigg].
\end{aligned}
\]

最后一项的期望为零。再利用因果性，第 \(t\) 个 token 的 score 不需要乘过去的 \(d_k,\;k<t\)，得到：

\[
\boxed{
\nabla_\theta\mathcal J
=
\mathbb E_{x\sim\pi_\theta}
\left[
\sum_{t=0}^{T-1}
\left(
\sum_{k=t}^{T-1}d_k
\right)
\nabla_\theta
\log\pi_\theta(x_{t+1}|x_{1..t})
\right].
}
\]

因此，严格的序列级反向 KL 梯度中，第 \(t\) 个 token 的权重应该是：

\[
d_t+d_{t+1}+\cdots+d_{T-1}.
\]

原因是当前 token 不仅影响当前 KL，还会改变后面访问到的所有前缀。

因此我们定义折扣未来 KL：

\[
G_t^{(\gamma)}
=
\sum_{k=t}^{T-1}
\gamma^{k-t}d_k.
\]

即：

\[
G_t^{(\gamma)}
=
d_t+\gamma d_{t+1}+\gamma^2d_{t+2}+\cdots.
\]

对应的 KL advantage 为：

\[
A_t^{(\gamma)}
=
-\beta G_t^{(\gamma)}.
\]

不同 \(\gamma\) 的含义如下。

\subsection{\texorpdfstring{\(\gamma=1\)}{gamma=1}：完整序列信用分配}

\[
G_t^{(1)}
=
d_t+d_{t+1}+\cdots+d_{T-1}.
\]

这对应序列级目标下的无折扣信用分配。

\subsection{\texorpdfstring{\(0<\gamma<1\)}{0<gamma<1}：部分未来信用分配}

\[
G_t^{(\gamma)}
=
d_t+\gamma d_{t+1}+\gamma^2d_{t+2}+\cdots.
\]

保留未来影响，但越远的 KL 权重越小。

\subsection{\texorpdfstring{\(\gamma=0\)}{gamma=0}：只看即时 token}

\[
\boxed{
G_t^{(0)}=d_t.
}
\]

当前 token 只根据当前位置的教师信号更新，不承担未来 KL。这也是目前的默认选择。主要是因为 TML 官方博客里有提及： 正折扣因子在数学上更完整，但实验中没有观察到性能改善，因此为了简单，主动使用零折扣。

\section{OPD 的最终统一公式}

\subsection{Rollout}

\[
x_{t+1}
\sim
q(\cdot|x_{1..t}).
\]

\subsection{sampled reverse KL}

\[
d_t
=
\log q(x_{t+1}|x_{1..t})
-
\log\pi_{\text{teacher}}(x_{t+1}|x_{1..t}).
\]

\subsection{折扣未来 KL}

\[
G_t^{(\gamma)}
=
\sum_{k=t}^{T-1}
\gamma^{k-t}d_k.
\]

\subsection{advantage}

\[
\boxed{
A_t
=
-\beta\operatorname{stopgrad}
\left[
G_t^{(\gamma)}
\right].
}
\]

\subsection{importance ratio}

\[
\boxed{
\rho_t(\theta)
=
\frac{
p_\theta(x_{t+1}|x_{1..t})
}{
q(x_{t+1}|x_{1..t})
}.
}
\]

\subsection{Tinker loss}

\[
\boxed{
\mathcal L_{\mathrm{OPD}}
=
-\sum_t
\rho_t(\theta)A_t.
}
\]

\subsection{梯度}

\[
\boxed{
\nabla_\theta\mathcal L_{\mathrm{OPD}}
=
\beta
\sum_t
\rho_t(\theta)
G_t^{(\gamma)}
\nabla_\theta
\log p_\theta(x_{t+1}|x_{1..t}).
}
\]

默认：

\[
\gamma=0,
\]

所以：

\[
G_t^{(0)}=d_t,
\]

最终为：

\[
\boxed{
\nabla_\theta\mathcal L_{\mathrm{OPD}}
=
\beta
\sum_t
\frac{p_\theta(x_{t+1}|x_{1..t})}
{q(x_{t+1}|x_{1..t})}
\left[
\log q(x_{t+1}|x_{1..t})
-
\log\pi_{\text{teacher}}(x_{t+1}|x_{1..t})
\right]
\nabla_\theta
\log p_\theta(x_{t+1}|x_{1..t}).
}
\]

严格 on-policy 时：

\[
p_\theta=q=\pi_\theta,
\]

化简为：

\[
\boxed{
\nabla_\theta\mathcal L_{\mathrm{OPD}}
=
\beta
\sum_t
\left[
\log\pi_\theta(x_{t+1}|x_{1..t})
-
\log\pi_{\text{teacher}}(x_{t+1}|x_{1..t})
\right]
\nabla_\theta
\log\pi_\theta(x_{t+1}|x_{1..t}).
}
\]




\end{document}
